\section{Apparato Bibliografico e Nota Metodologica}

\subsection{Dichiarazione sull'uso di Strumenti di Intelligenza Artificiale Generativa}

Nel corso della redazione del presente elaborato è stato impiegato uno strumento di Intelligenza Artificiale Generativa a supporto delle attività di analisi, revisione e organizzazione dei contenuti.

L'utilizzo dell'AI non ha sostituito il contributo intellettuale del gruppo, ma ha rappresentato un supporto metodologico nelle seguenti attività.

\subsubsection*{1. Supporto alla scelta del dominio applicativo}

In una fase preliminare è stato richiesto all'AI un confronto sintetico tra diversi domini IoT (Domotica, Smart City, Medical, Industrial, Agricoltura, Logistica, Fleet Monitoring), al fine di evidenziare le principali criticità ingegneristiche tipiche di ciascun settore.

In particolare, sono stati analizzati aspetti quali:
\begin{itemize}
    \item requisiti di latenza (real-time critico vs best-effort),
    \item vincoli energetici,
    \item problematiche di sicurezza e privacy,
    \item complessità dell'ambiente operativo,
    \item scalabilità del sistema.
\end{itemize}

Questo confronto ha permesso al gruppo di maturare una scelta consapevole del dominio ritenuto soggettivamente più interessante, non in termini di semplicità, ma di ricchezza di trade-off progettuali.

La decisione finale è stata presa autonomamente dal gruppo, sulla base di considerazioni critiche.

\subsubsection*{2. Revisione sintattica e miglioramento espositivo}

L'AI è stata utilizzata per effettuare revisioni linguistiche su testi già redatti dal gruppo, con l'obiettivo di:
\begin{itemize}
    \item migliorare la chiarezza espositiva,
    \item rendere più rigoroso il linguaggio tecnico,
    \item evitare ripetizioni e ambiguità sintattiche.
\end{itemize}

I contenuti tecnici, le analisi dei requisiti e le valutazioni ingegneristiche restano frutto esclusivo dell'elaborazione del gruppo.

\subsubsection*{3. Supporto alla fase esplorativa del caso di studio}

Per il Capitolo 3, relativo all'analisi di un caso di studio reale, è stato richiesto all'AI di individuare diversi possibili casi di applicazioni IoT coerenti con il dominio selezionato.

Lo strumento è stato utilizzato come supporto esplorativo con lo scopo di:
\begin{itemize}
    \item ridurre il tempo necessario alla fase preliminare di ricerca,
    \item filtrare i casi meno pertinenti,
    \item concentrare l'attenzione su soluzioni effettivamente implementate nel mondo reale.
\end{itemize}

La selezione finale del caso di studio è stata effettuata dal gruppo, sulla base della coerenza con i requisiti progettuali identificati nel Capitolo 2.

\subsubsection*{4. Sintesi preliminare di articoli e paper scientifici}

L'AI è stata utilizzata per ottenere riassunti preliminari di articoli scientifici e white paper individuati durante la ricerca bibliografica.

Tali sintesi sono state impiegate esclusivamente come:
\begin{itemize}
    \item strumento di orientamento iniziale,
    \item supporto alla comprensione generale del documento.
\end{itemize}

La lettura critica delle fonti, l'interpretazione dei risultati e l'integrazione nel presente elaborato sono state svolte direttamente dal gruppo.

\subsubsection*{5. Supporto alla formattazione \LaTeX}

Lo strumento è stato utilizzato per:
\begin{itemize}
    \item organizzare alcune sezioni secondo una struttura coerente,
    \item verificare la correttezza sintattica di parti in \LaTeX,
    \item uniformare l'impostazione delle sezioni e sottosezioni.
\end{itemize}

Non è stato effettuato alcun copia-incolla automatico di contenuti generati senza revisione critica.

\subsection{Dichiarazione del Contributo Originale}

Il contributo originale del gruppo si è concretizzato principalmente nelle seguenti attività:
\begin{itemize}
    \item Analisi critica dei requisiti di connettività, latenza, energia e sicurezza nel dominio selezionato.
    \item Identificazione dei trade-off ingegneristici (energia vs banda, latenza vs consumo, sicurezza vs overhead computazionale).
    \item Valutazione della coerenza tra il caso di studio analizzato e i modelli architetturali IoT discussi a lezione (architetture a 3, 5 e 7 livelli).
    \item Sintesi dei vincoli progettuali che guideranno la successiva fase di design dell'architettura.
\end{itemize}

L'AI ha svolto un ruolo di supporto metodologico, ma l'elaborazione concettuale, l'analisi tecnica e le conclusioni rappresentano un apporto intellettuale esclusivo del gruppo.